\section{Related Work}

\noindent \textbf{Indoor Localization via Wireless Networks}
Indoor localization based on existing wireless communication networks, such as WiFi \cite{wifireview2015yang} and Bluetooth\cite{bluetoo2016chen}, based on their RSS, has gained popularity because it does not require additional hardware \cite{wifireview2015yang}. For such systems, two primary methodologies are employed:
1) Utilizing an access-point database that comprises access-point IDs and their respective geo-coordinates. This setup facilitates the determination of user location through trilateration or weighted averaging of access-point coordinates \cite{WiGLEonline}.
2) Employing a fingerprint database, a radiomap, which records the received signal strength indicators and the corresponding geo-coordinates at their observation points\cite{radar2000bah},\cite{WILL2013wu}. Such an approach stands out for its inherent scalability, amplified by the prospects of crowd-sourced data collection \cite{crowdsource2010mink}. 

% \textbf{TBD...}\chris{@Josh, do you still put it a TBD?}

\noindent \textbf{Uncertainty Estimation}
Estimating the uncertainty, or more specifically, the suspicious evidence cues derived from data is crucial for the integration in real-world applications and as input for data fusion processes. 
To quantify uncertainty within deep learning, two primary distinctions emerge: epistemic uncertainty, associated with the model, and aleatoric uncertainty, related to the data \cite{kendall2017uncertainties}. 
Estimating aleatoric uncertainty can enhance the robustness of the network when training on noisy datasets. 
By introducing an additional aleatoric uncertainty prediction head, the reliability of neural networks has been significantly strengthened across various applications \cite{cai2022stun}\cite{poggi2020uncertainty}\cite{taha2019unsupervised}\cite{feng2018safe}. Our research primarily concentrates on enhancing localization performance on noisy radio maps by predicting uncertainty, which in turn aids in data fusion processes. 


\noindent \textbf{Multi-modal Sensor Fusion}
Indoor positioning and navigation techniques using smartphones have seen significant advances with the incorporation of various data sources and modeling methods. WiFi-SLAM \cite{ferris2007wifi,mirowski2013signalslam} relays on WiFi-based loop-closing constraints, which effectively minimize the accumulation errors commonly found in inertial navigation. Bayesian filter-based algorithms can also be used to perform the multi-modal sensor fusion to estimate the online position, such as particle filter\cite{zou2017Accurate,Hong2014WaP,carrera2016real}, Kalman filter and its variant\cite{Hellmers2013IMU,Poulose2019Indoor,Feng2020Kalman}. Several studies have incorporated floor plan information into their fusion systems via offline image processing techniques \cite{herath2021fusiondhl,ma2017pedestrian}, while other methods use the floor plan as a constraint for particle filtering \cite{Hong2014WaP,carrera2016real} and conditional random fields\cite{xiao2014lightweight}. However, there is no method to properly fuse IMU, uncertainty-aware WiFi localization results, and the crowdsourced WiFi radio map data.
